\documentclass[12pt]{article}
\usepackage[utf8]{inputenc}
\usepackage[T1]{fontenc}
\usepackage{hyperref}
\title{Ukuriguma Language Resources for K-12 Curriculum Development}
\author{Compiled by Claude Sonnet 5 (claude-sonnet-5) for Andr\'as Kornai}
\date{}
\begin{document}
\maketitle
\section{Introduction}
Ukuriguma (ISO 639-3 ukg) is a Trans-New Guinea language of the North Adelbert group, spoken in Madang Province, Papua New Guinea, roughly 40 km northwest of Madang town. Wikipedia gives 170 speakers (2003), and Joshua Project estimates about 300 people. The Endangered Languages Project lists it as endangered.
\section{Primary Educational Resources}
No educational materials were found in the search results.
\section{Linguistic and Reference Materials}
\begin{itemize}
\item Wikipedia, Ukuriguma language: \url{https://en.wikipedia.org/wiki/Ukuriguma_language}
\item ASJP word list entry: \url{https://asjp.clld.org/languages/UKURIGUMA}
\item Z'graggen 1980 Northern Adelbert word list, via Trans New Guinea database: \url{https://transnewguinea.org/source/zgraggen-1980-nadelbert?page=79}
\end{itemize}
\section{Religious and Community Resources}
\begin{itemize}
\item Joshua Project profile: \url{https://joshuaproject.net/people_groups/15698/PP}
\end{itemize}
\section{Digital and Audio Resources}
\begin{itemize}
\item Endangered Languages Project entry: \url{https://www.endangeredlanguages.com/lang/3433}
\end{itemize}
\section{Curriculum Development Context}
Only word lists and basic reference entries exist. A curriculum would require new documentation work with the community.
\end{document}